\NeedRoom{10}

\section{Data, Assumptions, and Notation}

The recovered TransLink feed is 23SEP\_20231020, declared for 4 September-31 December 2023 \cite{ref8}. The adapter uses the existing cleaned trip manifest, calendar and exceptions. Times beyond 24:00 remain on the service-day clock. Stop IDs are not merged by name.


\begin{longtable}{@{}L{\dimexpr 0.23000\textwidth-2\tabcolsep\relax}L{\dimexpr 0.37000\textwidth-2\tabcolsep\relax}L{\dimexpr 0.40000\textwidth-2\tabcolsep\relax}@{}}
\caption{Scheduling inputs, assumptions and evidence status}\label{tab:source2}\\
\toprule
\textbf{Input} & \textbf{Central treatment} & \textbf{Evidence status} \\
\midrule
\endfirsthead
\multicolumn{3}{l}{Table \thetable\ (continued)}\\
\toprule
\textbf{Input} & \textbf{Central treatment} & \textbf{Evidence status} \\
\midrule
\endhead
\midrule\multicolumn{3}{r}{Continued on next page}\\\endfoot
\bottomrule\endlastfoot
Times / endpoints & Preserved trip obligations & Published schedule, not actual operations \\
Trip distance & Audited upper interval endpoint & Derived bound; lower endpoint tested \\
Travel matrix & 1.35 x great-circle km; 30 km/h & Assumed road detour and speed \\
Candidate sites & 534 endpoint IDs & Coordinates observed; buildability unknown \\
Trolley service & Retain cap47/cap70 proxy separately & Existing proxy, not propulsion roster \\
Capacity / vehicle types & Homogeneous, equal capacity & No individual load or axle verification \\
R / T & 270 km / 60 min & Chosen usable range and charge duration \\
Energy / efficiency & 1.60 kWh/km / 90\% & Historical all-electric transfer \cite{ref9} / assumption \\
Night / depot & Full at first trip; no pullout or pullin & Single-day boundary; cross-day closure absent \\
Charging capacity & Unlimited in solver; count slots afterward & No utility or site certificate \\
\end{longtable}

The travel-time matrix is metric. This ensures a route via a charger is not faster than the direct transfer, which is necessary for the diesel lower-bound proof. A regression test shows that arbitrary nonmetric inputs can violate that bound.


All 25 distinct service patterns are evaluated separately. We use their maximum fleet requirement and retain their frequencies over 119 days. Annual activity is multiplied by 365/119. It is a seasonal extrapolation, not observed annual mileage.


\begin{align}
B_{\mathrm{use}}&=\varepsilon R, & B&=B_{\mathrm{use}}/u,\\
P_{\mathrm{battery}}&=60B_{\mathrm{use}}/T, &
P_{\mathrm{grid}}&=P_{\mathrm{battery}}/\eta .
\end{align}

With u = 0.8, the central battery is 540 kWh and grid-side charging power is 480 kW. The factor 60 corrects the minute/hour mismatch. A partial battery still occupies the charger for T, but receives only the energy consumed since its last full charge.


\subsection*{Principal notation}
\begin{longtable}{@{}L{0.22\textwidth}L{0.70\textwidth}@{}}
\caption{Principal notation and units}\label{tab:notation}\\
\toprule Symbol & Meaning and units\\\midrule
$i,j;\ n$ & Trip indices; total scheduled trips.\\
$o_i,d_i;\ s_i,e_i$ & Origin and destination; departure and arrival (minutes).\\
$\ell_i;\ u_i$ & Service distance; distance since the last full charge (km).\\
$w(a,b);\ \tau(a,b)$ & Deadhead distance (km) and travel time (minutes).\\
$R;\ T$ & Usable range (km); fixed charge duration (minutes).\\
$S;\ K$ & Selected charging nodes; node-count limit.\\
$N_d,N_e;\ \mathrm{PVR}$ & Diesel minimum, feasible electric fleet, concurrent-trip lower bound.\\
$\varepsilon;\ \eta;\ u$ & Battery kWh/km; charging efficiency; usable battery fraction.\\
$\gamma_t;\ \phi$ & Grid kgCO$_2$e/kWh; eligible-capital grant share.\\
\bottomrule
\end{longtable}
The state $u_i$ and the battery fraction $u$ are distinct quantities.
